% GENERATED FILE. Edit canonical lessons, then run npm run generate:books.
% canonical-source-hash: fnv1a64:32c0e527114d2935
% canonical-lessons: JA-C227-machigaeru, JA-C227-kangaeru, JA-C227-omoidasu, JA-C227-tashikameru, JA-C227-kuraberu

\chapter{Mistakes and Second Thoughts}
\label{ch:ja-mistakes-and-second-thoughts}
\hlchaptermodality{227}

\begin{chapteropening}
\textbf{By the end of this chapter:} \emph{I can talk about mistakes, thinking and checking.}

Complete the last lesson of chapter 227: I can talk about mistakes, thinking and checking.
\end{chapteropening}

\section[\texorpdfstring{\ja{まちがえる} (machigaeru)}{まちがえる (machigaeru)}]{\ja{まちがえる} (machigaeru) — to make a mistake; to get wrong}

\label{lesson:JA-C227-machigaeru}



\begin{quote}
Before the new one: say the Japanese for to take an exam, then the Japanese for a mark.
\end{quote}

\subsection*{You'll want to know: \ja{まちがえる}}
\textbf{\ja{まちがえる}} — \emph{machigaeru} — \textquotedblleft{}to make a mistake; to get wrong\textquotedblright{}.

\begin{quote}
\textbf{\ja{おとうとは} \ja{よく} \ja{なまえを} \ja{まちがえます。}} — \emph{otōto wa yoku namae o machigaemasu} — \textquotedblleft{}My younger brother often gets names wrong.\textquotedblright{}
\end{quote}

Polite \textbf{\ja{まちがえます}}, \emph{machigaemasu}; an ichidan verb.

\begin{cousinweb}
\textbf{\ja{まちがえる}} (\emph{machigaeru}) is a native Japanese word (wago), inherited rather than borrowed: it is built from \emph{ma}, \textquotedblleft{}interval, space\textquotedblright{}, and \emph{chigaeru}, the partner of \emph{chigau}, \textquotedblleft{}to differ\textquotedblright{} (JMdict, KANJIDIC2).
\end{cousinweb}

\begin{grammarlens}[title={What you've built}]
One more word for everyday life at A2.
\end{grammarlens}

\subsection*{Your turn}
\begin{itemize}
\raggedright
  \item \emph{Say it:} \emph{machigaeru}
  \item \emph{Say it:} \emph{machigaeru}, once more
  \item \emph{Say it:} \emph{otōto wa yoku namae o machigaemasu} — My younger brother often gets names wrong.
  \item \emph{Recall:} say the Japanese for to take an exam, then the Japanese for a mark, then say \emph{machigaeru} again
\end{itemize}

\begin{tcolorbox}[breakable,colback=teal!4,colframe=teal!35!black,title={Before you move on}]
What does \ja{まちがえる} mean? (\textquotedblleft{}To make a mistake; to get wrong\textquotedblright{}.) Say it once more.
\end{tcolorbox}
\section[\texorpdfstring{\ja{かんがえる} (kangaeru)}{かんがえる (kangaeru)}]{\ja{かんがえる} (kangaeru) — to think over; to consider}

\label{lesson:JA-C227-kangaeru}



\begin{quote}
Before the new one: say the Japanese for a mark, then the Japanese for to make a mistake.
\end{quote}

\subsection*{You'll want to know: \ja{かんがえる}}
\textbf{\ja{かんがえる}} — \emph{kangaeru} — \textquotedblleft{}to think over; to consider\textquotedblright{}.

\begin{quote}
\textbf{\ja{ちちは} \ja{いつも} \ja{ゆっくり} \ja{かんがえます。}} — \emph{chichi wa itsumo yukkuri kangaemasu} — \textquotedblleft{}My father always thinks things over slowly.\textquotedblright{}
\end{quote}

Polite \textbf{\ja{かんがえます}}, \emph{kangaemasu}; an ichidan verb.

\begin{cousinweb}
\textbf{\ja{かんがえる}} (\emph{kangaeru}) is a native Japanese word (wago), inherited rather than borrowed from Chinese; its kanji was attached later as a kun'yomi (JMdict, KANJIDIC2).
\end{cousinweb}

\begin{grammarlens}[title={What you've built}]
One more word for everyday life at A2.
\end{grammarlens}

\subsection*{Your turn}
\begin{itemize}
\raggedright
  \item \emph{Say it:} \emph{kangaeru}
  \item \emph{Say it:} \emph{kangaeru}, once more
  \item \emph{Say it:} \emph{chichi wa itsumo yukkuri kangaemasu} — My father always thinks things over slowly.
  \item \emph{Recall:} say the Japanese for a mark, then the Japanese for to make a mistake, then say \emph{kangaeru} again
\end{itemize}

\begin{tcolorbox}[breakable,colback=teal!4,colframe=teal!35!black,title={Before you move on}]
What does \ja{かんがえる} mean? (\textquotedblleft{}To think over; to consider\textquotedblright{}.) Say it once more.
\end{tcolorbox}
\section[\texorpdfstring{\ja{おもいだす} (omoidasu)}{おもいだす (omoidasu)}]{\ja{おもいだす} (omoidasu) — to recall; to remember}

\label{lesson:JA-C227-omoidasu}



\begin{quote}
Before the new one: say the Japanese for to make a mistake, then the Japanese for to think over.
\end{quote}

\subsection*{You'll want to know: \ja{おもいだす}}
\textbf{\ja{おもいだす}} — \emph{omoidasu} — \textquotedblleft{}to recall; to remember\textquotedblright{}.

\begin{quote}
\textbf{\ja{ははは} \ja{よく} \ja{なつやすみを} \ja{おもいだします。}} — \emph{haha wa yoku natsuyasumi o omoidashimasu} — \textquotedblleft{}My mother often remembers the summer holidays.\textquotedblright{}
\end{quote}

Polite \textbf{\ja{おもいだします}}, \emph{omoidashimasu}; a godan verb.

\begin{cousinweb}
\textbf{\ja{おもいだす}} (\emph{omoidasu}) is a native Japanese compound (wago), inherited rather than borrowed: it is built from \emph{omoi}, from \emph{omou}, \textquotedblleft{}to think\textquotedblright{}, and \emph{dasu}, \textquotedblleft{}to take out\textquotedblright{} (JMdict, KANJIDIC2).
\end{cousinweb}

\begin{grammarlens}[title={What you've built}]
One more word for everyday life at A2.
\end{grammarlens}

\subsection*{Your turn}
\begin{itemize}
\raggedright
  \item \emph{Say it:} \emph{omoidasu}
  \item \emph{Say it:} \emph{omoidasu}, once more
  \item \emph{Say it:} \emph{haha wa yoku natsuyasumi o omoidashimasu} — My mother often remembers the summer holidays.
  \item \emph{Recall:} say the Japanese for to make a mistake, then the Japanese for to think over, then say \emph{omoidasu} again
\end{itemize}

\begin{tcolorbox}[breakable,colback=teal!4,colframe=teal!35!black,title={Before you move on}]
What does \ja{おもいだす} mean? (\textquotedblleft{}To recall; to remember\textquotedblright{}.) Say it once more.
\end{tcolorbox}
\section[\texorpdfstring{\ja{たしかめる} (tashikameru)}{たしかめる (tashikameru)}]{\ja{たしかめる} (tashikameru) — to make sure; to check carefully}

\label{lesson:JA-C227-tashikameru}



\begin{quote}
Before the new one: say the Japanese for to think over, then the Japanese for to recall.
\end{quote}

\subsection*{You'll want to know: \ja{たしかめる}}
\textbf{\ja{たしかめる}} — \emph{tashikameru} — \textquotedblleft{}to make sure; to check carefully\textquotedblright{}.

\begin{quote}
\textbf{\ja{がくせいは} \ja{こたえを} \ja{もういちど} \ja{たしかめます。}} — \emph{gakusei wa kotae o mō ichido tashikamemasu} — \textquotedblleft{}The students check their answers once more.\textquotedblright{}
\end{quote}

Polite \textbf{\ja{たしかめます}}, \emph{tashikamemasu}; an ichidan verb.

\begin{cousinweb}
\textbf{\ja{たしかめる}} (\emph{tashikameru}) is a native Japanese word (wago), inherited rather than borrowed: it is built on \emph{tashika}, \textquotedblleft{}certain\textquotedblright{}, with a verb ending (JMdict, KANJIDIC2).
\end{cousinweb}

\begin{grammarlens}[title={What you've built}]
One more word for everyday life at A2.
\end{grammarlens}

\subsection*{Your turn}
\begin{itemize}
\raggedright
  \item \emph{Say it:} \emph{tashikameru}
  \item \emph{Say it:} \emph{tashikameru}, once more
  \item \emph{Say it:} \emph{gakusei wa kotae o mō ichido tashikamemasu} — The students check their answers once more.
  \item \emph{Recall:} say the Japanese for to think over, then the Japanese for to recall, then say \emph{tashikameru} again
\end{itemize}

\begin{tcolorbox}[breakable,colback=teal!4,colframe=teal!35!black,title={Before you move on}]
What does \ja{たしかめる} mean? (\textquotedblleft{}To make sure; to check carefully\textquotedblright{}.) Say it once more.
\end{tcolorbox}
\section[\texorpdfstring{\ja{くらべる} (kuraberu)}{くらべる (kuraberu)}]{\ja{くらべる} (kuraberu) — to compare}

\label{lesson:JA-C227-kuraberu}



\begin{quote}
Before the new one: say the Japanese for to recall, then the Japanese for to make sure.
\end{quote}

\subsection*{You'll want to know: \ja{くらべる}}
\textbf{\ja{くらべる}} — \emph{kuraberu} — \textquotedblleft{}to compare\textquotedblright{}.

\begin{quote}
\textbf{\ja{がくせいは} \ja{ふたつの} \ja{こたえを} \ja{くらべます。}} — \emph{gakusei wa futatsu no kotae o kurabemasu} — \textquotedblleft{}The students compare two answers.\textquotedblright{}
\end{quote}

Polite \textbf{\ja{くらべます}}, \emph{kurabemasu}; an ichidan verb.

\begin{cousinweb}
\textbf{\ja{くらべる}} (\emph{kuraberu}) is a native Japanese word (wago), inherited rather than borrowed from Chinese; its kanji was attached later as a kun'yomi (JMdict, KANJIDIC2).
\end{cousinweb}

\begin{grammarlens}[title={What you've built}]
One more word for everyday life at A2.
\end{grammarlens}

\subsection*{Your turn}
\begin{itemize}
\raggedright
  \item \emph{Say it:} \emph{kuraberu}
  \item \emph{Say it:} \emph{kuraberu}, once more
  \item \emph{Say it:} \emph{gakusei wa futatsu no kotae o kurabemasu} — The students compare two answers.
  \item \emph{Recall:} say the Japanese for to recall, then the Japanese for to make sure, then say \emph{kuraberu} again
\end{itemize}

\begin{tcolorbox}[breakable,colback=teal!4,colframe=teal!35!black,title={Before you move on}]
What does \ja{くらべる} mean? (\textquotedblleft{}To compare\textquotedblright{}.) Say it once more.
\end{tcolorbox}
